% FORRT LaTeX Curriculum Vitae Template
% Author: Emily Friedel
% Website: https://experts.deakin.edu.au/65529-emily-friedel
% Last updated: February 2026

% You may use use this document as a template to create your own CV
% and you may redistribute the source code freely. No attribution is
% required in any resulting documents. Please leave this notice and 
% the above URL in the source code if you choose to redistribute this file.

\documentclass{article}

\title{CV}
\date{February 2026}

% Packages
\usepackage{hyperref}
\usepackage{geometry}
\usepackage{graphicx}
\usepackage[T1]{fontenc}
\usepackage[sc]{mathpazo}
\usepackage{enumitem}
\usepackage[super]{nth}
\usepackage[dvipsnames]{xcolor}
\usepackage{multicol}
\usepackage{setspace}

% Page layout
\geometry{
    body={7in, 9.0in},
    left=0.7in,
    top=0.75in
}

\setstretch{1.12}

% Hyperlink setup
\hypersetup{
    colorlinks=true,
    linkcolor=Blue,
    urlcolor=Blue,
    citecolor=Blue,
    filecolor=Blue
}

% Name and contact information
\def\name{Yoshitaka Inoue}
\def\position{Ph.D. Candidate in Computer Science \,$\cdot$\, NIH Predoctoral Fellow}
\def\affiliation{University of Minnesota, Twin Cities \,$\cdot$\, National Institutes of Health}
\def\address{Bethesda, Maryland}
\def\email{\href{mailto:inoue.yoshitaka.iv1@gmail.com}{inoue.yoshitaka.iv1@gmail.com}}
\def\website{\href{https://inoue0426.github.io/}{inoue0426.github.io}}
\def\scholar{\href{https://scholar.google.co.jp/citations?user=aBizyLkAAAAJ\&hl=en}{Google Scholar}}
\def\github{\href{https://github.com/inoue0426}{GitHub}}
\def\linkedin{\href{https://www.linkedin.com/in/inoue0426/}{LinkedIn}}

\begin{document}

% Header
\begin{center}
    {\huge \textbf{\name}}\\
    \vspace{0.2em}
    {\large \position}\\
    \affiliation\\
    \address\\
    \email \, $\cdot$ \, \website\\
    \scholar \, $\cdot$ \, \github \, $\cdot$ \, \linkedin
\end{center}

% Section: Education
\section*{Education}
\hrule
\vspace{1em}
\begin{itemize}[leftmargin=*]
    \item[] \textbf{Ph.D. Candidate in Computer Science and Engineering}, University of Minnesota, Twin Cities, USA \hfill 2021 -- Present\\
    Advisor: Prof. Rui Kuang; Co-advisor: Dr. Augustin Luna (NIH)\\
    Expected graduation: August 2027
    
    \item[] \textbf{M.Eng. in Information Science}, Nara Institute of Science and Technology, Japan \hfill 2017 -- 2019\\
    Advisor: Prof. Shigehiko Kanaya
    
    \item[] \textbf{B.Sc. in Biological Science}, Kumamoto University, Japan \hfill 2012 -- 2017\\
    Advisor: Prof. Susumu Takio
\end{itemize}

% Section: Research Interests
\section*{Research Interests}
\hrule
\vspace{1em}
\begin{itemize}[leftmargin=2em]
    \item Machine learning for biological perturbations and therapeutic response
    \item Single-cell and functional genomics
    \item Intervention-aware representation learning
    \item Graph-based biomedical reasoning and LLM-based scientific reasoning
\end{itemize}

% Section: Research Experience
\section*{Research Experience}
\hrule
\vspace{1em}
\begin{itemize}[leftmargin=*]
    \item[] \textbf{Pre-doctoral Fellow}, National Library of Medicine / National Cancer Institute, NIH \hfill 2024 -- Present\\
    Co-advisor: Dr. Augustin Luna
    \begin{itemize}
        \item Develop machine learning methods for biological perturbation modeling and therapeutic response prediction across single-cell, pharmacogenomic, high-throughput screening, cancer cell-line, and patient datasets.
        \item Build treatment-conditioned latent-state transition models for transferring perturbation-response information to patient drug-response prediction.
        \item Develop graph-based drug-response models and biomedical reasoning systems that integrate molecular, pharmacologic, and clinical evidence.
    \end{itemize}

    \item[] \textbf{Google Summer of Code Contributor; Mentor}, Google Summer of Code \hfill 2022 -- 2024
    \begin{itemize}
        \item Built a drug-response prediction pipeline and mentored projects on PyTorch Geometric datasets and multi-tool LLM aggregation for biomedical analysis.
    \end{itemize}

    \item[] \textbf{Machine Learning Engineer}, FreakOut, Inc., Tokyo, Japan \hfill 2019 -- 2021
    \begin{itemize}
        \item Applied machine learning to click-through-rate prediction.
    \end{itemize}
\end{itemize}

% Section: Peer-Reviewed Publications
\section*{Peer-Reviewed Publications}
\hrule
\vspace{1em}
\begin{itemize}[leftmargin=*]
    \item[] \textbf{Inoue, Y.}, Lee, H., Fu, T., Kuang, R., \& Luna, A. (2026).
    drGT: Attention-Guided Gene Assessment of Drug Response Utilizing a Drug-Cell-Gene Heterogeneous Network.
    \textit{BMC Bioinformatics}.

    \item[] \textbf{Inoue, Y.}, Hao, N., Lu, Y., Fu, T., \& Luna, A. (2026).
    Drug Discovery in the Era of Artificial Intelligence: From Target Identification to Clinical Trials.
    \textit{Health Data Science}. 

    \item[] Hoshi-Kadonosawa, Y., Reinhold, W. C., Taniyama, D., \textbf{Inoue, Y.}, et al. (2026).
    TRPM4 Expression as a Predictive Biomarker and a Mechanistic Driver of Acetalax Activity in Prostate Cancer: Preclinical Efficacy Studies.
    \textit{Molecular Cancer Therapeutics}.

    \item[] Elloumi, F., Reinhold, W. C., Varma, S., Wang, Y., Kinali, M., Arakawa, Y., \textbf{Inoue, Y.}, et al. (2026).
    CellMiner cross-database (CellMinerCDB) version 2.2 for explorations of patient-derived cancer cell line pharmacogenomics.
    \textit{Nucleic Acids Research}.
\end{itemize}

% Section: Preprints and Manuscripts in Preparation
\section*{Preprints and Manuscripts Under Review}
\hrule
\vspace{1em}
\begin{itemize}[leftmargin=*]
    \item[] \textbf{Inoue, Y.}, Jeong, M., Hero, A., Kuang, R., \& Luna, A. (2026).
    PerturbRx: Learning Treatment-Conditioned Latent Transitions for Patient Drug Response Prediction.
    \textit{arXiv}.

    \item[] \textbf{Inoue, Y.}, Song, T., Wang, X., Luna, A., \& Fu, T. (2026).
    DrugAgent: Reliable Multi-Agent Aggregation under Conflicting Biomedical Evidence.
    \textit{Scientific Reports}. Under revision.

    \item[] Taniyama, D., \textbf{Inoue, Y.}, Elloumi, F., Kim, Y. S., et al. (2026).
    Spontaneous dynamic gene expression of Schlafen 11 (SLFN11) demonstrated by single-cell RNA sequencing in clonal cancer cell lines.
    \textit{Oncogene}. Submitted.

    \item[] \textbf{Inoue, Y.} (2024).
    scVGAE: A Novel Approach Using ZINB-Based Variational Graph Autoencoder for Single-Cell RNA-seq Imputation.
    \textit{arXiv}.

    \item[] \textbf{Inoue, Y.}, Kulman, E., \& Kuang, R. (2024).
    BiGCN: Leveraging Cell and Gene Similarities for Single-cell Transcriptome Imputation with Bi-Graph Convolutional Networks.
    \textit{bioRxiv}.
\end{itemize}

% Section: Conference and Workshop Papers
\section*{Conference and Workshop Papers}
\hrule
\vspace{1em}
\begin{itemize}[leftmargin=*]
    \item[] \textbf{Inoue, Y.}, Song, T., Wang, X., Luna, A., \& Fu, T. (2026).
    DrugAgent: Reliable Multi-Agent Aggregation under Conflicting Biomedical Evidence.
    \textit{ICML Workshop on Foundation Models for the Life Sciences (FM4LS)}.

    \item[] \textbf{Inoue, Y.}, Fu, T., \& Luna, A. (2025).
    GraphPINE: Graph Importance Propagation for Interpretable Drug Response Prediction.
    \textit{ICLR Workshop on Machine Learning for Genomics Explorations (MLGenX)}.

    \item[] \textbf{Inoue, Y.}, Song, T., Wang, X., Luna, A., \& Fu, T. (2025).
    DrugAgent: Explainable Drug Repurposing Agent with Large Language Model-based Reasoning.
    \textit{ICLR Workshop on Machine Learning for Genomics Explorations (MLGenX)}.
\end{itemize}

% Section: Conference Presentations
\section*{Selected Conference Presentations}
\hrule
\vspace{1em}
\begin{itemize}[leftmargin=*]
    \item[] \textbf{Inoue, Y.} (2026).
    drGT [Oral presentation].
    \textit{ISMB Network Biology (NetBio) Community of Special Interest}, Washington, DC.

    \item[] \textbf{Inoue, Y.} (2026).
    DrugAgent [Oral presentation].
    \textit{ISMB Text Mining Community of Special Interest}, Washington, DC.

    \item[] \textbf{Inoue, Y.} (2026).
    drGT [Oral presentation].
    \textit{26th Annual CCR-FYI Colloquium}, Rockville, MD.
    Selected as one of four Outstanding Postgraduate Fellow Finalists.

    \item[] \textbf{Inoue, Y.} (2025).
    drGT: attention-guided gene assessment for drug-response modeling [Oral presentation].
    \textit{NIH Artificial Intelligence Symposium}, Rockville, MD.
    Selected among 5 of 78 listed abstracts.
\end{itemize}

% % Section: Non-Traditional Research Outputs
% \section*{Research Software}
% \hrule
% \vspace{1em}
% \begin{itemize}[leftmargin=*]
%     \item[] \textbf{OncoGraph}, \href{https://github.com/inoue0426/OncoGraph}{github.com/inoue0426/OncoGraph}
%     \begin{itemize}
%         \item Open oncology knowledge graph and retrieval benchmark integrating drugs, genes, diseases, clinical trials, Gene Ontology terms, and provenance-linked biomedical evidence.
%     \end{itemize}
% \end{itemize}


% Section: Awards and Fellowships
\section*{Fellowships, Awards, and Honours}
\hrule
\vspace{1em}
\begin{itemize}[leftmargin=*]
    \item[] \textbf{NIH Predoctoral Research Fellowship} \hfill 2024--Present
    \item[] \textbf{Outstanding Postgraduate Fellow Finalist}, 26th Annual CCR-FYI Colloquium --- selected as 1 of 4 finalists \hfill 2026
    \item[] \textbf{Research Visit Scholarships}, NAIST and JASSO support for UC Davis research visits \hfill 2018--2019
\end{itemize}


% Section: Service roles
\section*{Academic Service}
\hrule
\vspace{1em}
\begin{itemize}[leftmargin=*]
    \item[] \textbf{Journal reviewing:}
    BMJ Digital Health \& AI (Expert Peer Reviewer, 2026--2027);
    IEEE/ACM Transactions on Computational Biology and Bioinformatics (2025--2027);
    Nature Computational Science (co-reviewer, 2026).

    \item[] \textbf{Conference reviewing:}
    ACM BCB (2023); ICML FM4LS (2026); ICML GenBio (2026); IEEE BIBM (2026); ACM SIGSPATIAL (2026).
\end{itemize}

% Section: Technical Skills
\section*{Technical Skills}
\hrule
\vspace{1em}
\begin{itemize}[leftmargin=*]
    \item[] \textbf{Machine Learning / AI:}
    Graph neural networks and heterogeneous graphs; autoencoders and latent-state models; representation learning; domain adaptation and transfer learning; treatment-conditioned transition modeling; LLM-based agents and retrieval-augmented generation; interpretable machine learning.

    \item[] \textbf{Computational Biology:}
    Single-cell and bulk RNA-seq; pharmacogenomics and drug-response modeling; perturbation and functional genomics; cancer genomics; biological network analysis; multi-omics integration; biomedical knowledge graphs.

    \item[] \textbf{Frameworks:}
    PyTorch; PyTorch Geometric; scikit-learn; LangChain; AutoGen.

    \item[] \textbf{Computational Infrastructure:}
    Slurm; HPC; GPU computing; CUDA; Docker; uv; Ollama; LightLLM.
\end{itemize}

% Footer
\begin{center}
    \begin{footnotesize}
        Last updated: \today
    \end{footnotesize}
\end{center}

\end{document}
